\section{Memória e Armazenamento de Valores}

A FX-82MS dispõe de recursos de memória que permitem armazenar e recuperar
valores intermediários durante a realização de cálculos. O uso adequado desses
recursos evita a necessidade de redigitar resultados obtidos em etapas anteriores,
reduzindo a possibilidade de erros de transcrição e tornando o processo de cálculo
mais eficiente. A calculadora oferece dois tipos principais de armazenamento: a
memória independente e as variáveis de A a F, M, X e Y \cite{casio2004}.

\subsection{Memória Independente}

A memória independente, identificada pela letra \textbf{M} no visor, permite
acumular valores ao longo de uma sequência de cálculos. Ela é gerenciada por
meio de três operações principais \cite{casio2004}:

\begin{itemize}
    \item \textbf{M+} -- adiciona o resultado atual ao valor armazenado em M;
    \item \textbf{SHIFT} seguido de \textbf{M+} (função \textbf{M$-$}) -- subtrai o resultado atual do valor armazenado em M;
    \item \textbf{RCL} seguido de \textbf{M+} -- recupera o valor armazenado em M e o exibe no visor.
\end{itemize}

O indicador \textbf{M} aparece no topo do visor sempre que há um valor diferente
de zero armazenado na memória independente. Para zerar a memória M sem apagar
as demais configurações, armazene o valor zero com a sequência
\textbf{0} $\to$ \textbf{M+}, ou utilize \textbf{SHIFT} $\to$ \textbf{CLR} $\to$
\textbf{1} (Mcl), que limpa todas as memórias simultaneamente
\cite{casio2004}.

Uma aplicação direta da memória independente é o cálculo de somas acumuladas,
como a soma dos quadrados dos desvios utilizada no cálculo da variância. Dadas as
notas 6, 7 e 8, com média $\bar{x} = 7$:

\begin{table}[H]
    \centering
    \begin{tabular}{lll}
        \hline
        \textbf{Etapa} & \textbf{Entrada} & \textbf{Resultado} \\ \hline
            Zerar a memória & \texttt{0 M+} & M = 0 \\
            Primeiro quadrado & \texttt{( 6 - 7 ) x\textasciicircum{}2 M+} & M = 1 \\
            Segundo quadrado & \texttt{( 7 - 7 ) x\textasciicircum{}2 M+} & M = 1 \\
            Terceiro quadrado & \texttt{( 8 - 7 ) x\textasciicircum{}2 M+} & M = 2 \\
            Recuperar o total & \texttt{RCL M+} & 2 \\
            Variância amostral & \texttt{Ans / ( 3 - 1 ) =} & 1 \\ \hline
        \end{tabular}
\end{table}

O valor acumulado em M permanece armazenado mesmo após pressionar \textbf{AC}
ou desligar a calculadora, o que exige atenção: recomenda-se sempre zerar a
memória independente antes de iniciar um novo cálculo que faça uso de M+,
evitando que valores residuais de operações anteriores interfiram no resultado.

\subsection{Variáveis de Armazenamento}

Além da memória independente, a FX-82MS disponibiliza oito variáveis de
armazenamento identificadas pelas letras \textbf{A}, \textbf{B}, \textbf{C},
\textbf{D}, \textbf{E}, \textbf{F}, \textbf{X} e \textbf{Y}, além da variável
\textbf{M}, compartilhada com a memória independente, totalizando nove variáveis
disponíveis \cite{casio2004}. Cada variável armazena um único valor
numérico e pode ser acessada individualmente, o que as torna adequadas para
guardar constantes ou resultados intermediários que serão reutilizados em
diferentes etapas de um mesmo cálculo \cite{casio2004}.

O armazenamento de um valor em uma variável é feito com a tecla \textbf{STO},
seguida da tecla correspondente à variável desejada. A recuperação do valor
armazenado é feita com a tecla \textbf{RCL}, também seguida da tecla da variável.
As letras A a F e M estão indicadas em vermelho acima de teclas específicas e são
acessadas por meio da tecla \textbf{ALPHA} \cite{casio2004}:

\begin{itemize}
    \item \textbf{STO} $\to$ \textbf{ALPHA} $\to$ \textit{tecla da variável} -- armazena o valor atual na variável indicada;
    \item \textbf{RCL} $\to$ \textbf{ALPHA} $\to$ \textit{tecla da variável} -- recupera o valor armazenado na variável indicada.
\end{itemize}

O conteúdo de cada variável permanece armazenado após pressionar \textbf{AC}
ou desligar a calculadora. Para limpar todas as variáveis simultaneamente,
utilize \textbf{SHIFT} $\to$ \textbf{CLR} $\to$ \textbf{1} (Mcl). A
reinicialização completa por meio de \textbf{SHIFT} $\to$ \textbf{CLR} $\to$
\textbf{3} (All) também apaga o conteúdo de todas as variáveis
\cite{casio2004}.

Um uso típico das variáveis é o armazenamento da média de um conjunto de dados
para reutilização no cálculo dos desvios individuais, evitando redigitar o valor
a cada etapa. Considere as notas 5, 7 e 9, cuja média é $\bar{x} = 7$:

\begin{table}[H]
    \centering
        \begin{tabular}{lll}
        \hline
        \textbf{Etapa} & \textbf{Entrada} & \textbf{Resultado} \\ \hline
            Calcular a média & \texttt{( 5 + 7 + 9 ) / 3 =} & 7 \\
            Armazenar em A & \texttt{STO ALPHA A} & 7 \\
            Desvio da nota 5 & \texttt{5 - RCL ALPHA A =} & -2 \\
            Desvio da nota 7 & \texttt{7 - RCL ALPHA A =} & 0 \\
            Desvio da nota 9 & \texttt{9 - RCL ALPHA A =} & 2 \\ \hline
        \end{tabular}
\end{table}

É importante distinguir o papel da variável \textbf{M} do papel da memória
independente, embora compartilhem o mesmo espaço de armazenamento: ao utilizar
\textbf{M+} ou \textbf{M$-$}, o valor em M é acumulado; ao utilizar
\textbf{STO} $\to$ \textbf{M}, o valor atual \emph{substitui} o conteúdo
anterior de M, sem acumulação. Essa distinção é relevante para evitar
comportamentos inesperados quando as duas formas de acesso à memória M são
combinadas em um mesmo cálculo \cite{casio2004}.

\subsection{Uso Estratégico em Cálculos Estatísticos}

A combinação entre a memória independente e as variáveis de armazenamento
permite estruturar cálculos estatísticos de forma organizada e sem retrabalho.
Uma estratégia eficiente consiste em armazenar valores fixos, como a média ou
um denominador comum, nas variáveis A a F, enquanto a memória M acumula
somas parciais ao longo das etapas \cite{casio2004}.

Essa abordagem é especialmente útil no cálculo da variância e do desvio padrão,
em que o mesmo valor da média é subtraído repetidamente de cada dado do
conjunto. Ao armazenar a média em uma variável antes de iniciar os cálculos dos
desvios, evita-se a redigitação do valor a cada etapa, reduzindo o risco de erro.

Ao final de uma sequência de cálculos que envolva M+, recomenda-se recuperar
o valor acumulado com \textbf{RCL M+} antes de prosseguir com as operações
seguintes. Quando o mesmo cálculo envolver tanto \textbf{STO M} quanto
\textbf{M+}, é importante ter clareza sobre qual operação está sendo utilizada
em cada etapa, uma vez que o comportamento de acumulação e de substituição
diferem, conforme descrito na seção anterior.

\subsection{Exemplos Práticos}

Os exemplos a seguir combinam o uso das variáveis de armazenamento e da
memória independente para resolver problemas completos.

\subsubsection*{Exemplo 1 -- Desvio padrão amostral}

Dado o conjunto de valores \{4, 7, 7, 10\}, calcule o desvio padrão amostral $s$. A média é $\bar{x} = (4 + 7 + 7 + 10)/4 = 7$. Armazena-se a média na variável A e acumulam-se os quadrados dos desvios na memória M:

\begin{table}[H]
    \centering
    \begin{tabular}{lll}
        \hline
        \textbf{Etapa} & \textbf{Entrada} & \textbf{Resultado} \\ \hline
            Calcular a média & \texttt{( 4 + 7 + 7 + 10 ) / 4 =} & 7 \\
            Armazenar em A & \texttt{STO ALPHA A} & 7 \\
            Zerar memória M & \texttt{0 M+} & M = 0 \\
            Desvio de 4 & \texttt{( 4 - RCL ALPHA A ) x\textasciicircum{}2 M+} & M = 9 \\
            Desvio de 7 (1º) & \texttt{( 7 - RCL ALPHA A ) x\textasciicircum{}2 M+} & M = 9 \\
            Desvio de 7 (2º) & \texttt{( 7 - RCL ALPHA A ) x\textasciicircum{}2 M+} & M = 9 \\
            Desvio de 10 & \texttt{( 10 - RCL ALPHA A ) x\textasciicircum{}2 M+} & M = 18 \\
            Recuperar soma & \texttt{RCL M+} & 18 \\
            Variância amostral & \texttt{Ans / ( 4 - 1 ) =} & 6 \\
            Desvio padrão & \texttt{$\sqrt{\phantom{0}}$ Ans =} & $\approx 2{,}449$ \\ \hline
    \end{tabular}
\end{table}

\subsubsection*{Exemplo 2 -- Média ponderada} 

Um aluno obteve as notas 6,0, 8,0 e 5,0 nas provas P1, P2 e P3, com pesos 2, 3 e 5, respectivamente. Calcule a média ponderada $\bar{x}$.

Armazena-se o denominador (soma dos pesos) na variável B e acumula-se o numerador (soma dos produtos nota $\times$ peso) na memória M:

\begin{table}[H]
    \centering
    \begin{tabular}{lll}
        \hline
        \textbf{Etapa} & \textbf{Entrada} & \textbf{Resultado} \\ \hline
            Armazenar soma dos pesos & \texttt{( 2 + 3 + 5 ) = STO ALPHA B} & 10 \\
            Zerar memória M & \texttt{0 M+} & M = 0 \\
            Produto P1 & \texttt{6 $\times$ 2 M+} & M = 12 \\
            Produto P2 & \texttt{8 $\times$ 3 M+} & M = 36 \\
            Produto P3 & \texttt{5 $\times$ 5 M+} & M = 61 \\
            Recuperar numerador & \texttt{RCL M+} & 61 \\
            Média ponderada & \texttt{Ans / RCL ALPHA B =} & 6{,}1 \\ \hline
    \end{tabular}
\end{table}